%% COMP7610 Technical Report shell — Potential Security Risk in ERC-721 Smart Contract
%% Based on IEEEtran peerreview template. Branch content: levi_dev lab docs.
\documentclass[peerreview]{IEEEtran}
\usepackage{cite}
\usepackage{url}
\usepackage[utf8]{inputenc}
\usepackage{booktabs}
\usepackage{graphicx}
\usepackage{hyperref}

\begin{document}

\title{Potential Security Risks in ERC-721 Smart Contracts:\\
Callback Reentrancy and Over-broad Approvals}

\author{Author One, Author Two, Author Three, Yifeng Liu\\
Department of Computer Science\\
The University of Hong Kong\\
COMP7610 Course Project}
\date{\today}

\maketitle
\tableofcontents
\listoffigures
\listoftables

\IEEEpeerreviewmaketitle

\begin{abstract}
% ~200 words. Cover: purpose; two threats (late app ledger on callback; setApprovalForAll over-breadth);
% method (Hardhat lab + public-report review); findings (CEI; least privilege); recommendations.
% Past tense / active voice for what was done. No exploit steps.
\end{abstract}

\section{Introduction}
% Motivation; two threat threads; report roadmap; defensive boundary (no exploit PoC).

\section{Problem Definition}
\subsection{Callback window and the two-ledger model}
% Nail three lines: (1) safe $\neq$ anti-reentrancy;
% (2) token ledger (owner/balance) vs app ledger (addressMinted);
% (3) same-tx callback window after \_safeMint.
% Cite lab: docs/CALLBACK-FLOW.md, docs/全流程-回调重入.md.

\subsection{Over-broad approvals}
% approve vs setApprovalForAll; public reports on BAYC \#3738 as approval phishing
% (not a BAYC contract bug; not callback reentrancy).

\subsection{Shared framing, distinct bugs}
% Handing over \emph{allowance} vs handing over \emph{control flow};
% motto: lock state before handing over power.

\section{Proposed Solutions}
\subsection{Application-level Checks-Effects-Interactions}
% Fixed: write addressMinted before \_safeMint; optional ReentrancyGuard as depth.

\subsection{Least-privilege approvals}
% Per-token approve; revocable operators; UI disclosure; no phishing playbook.

\subsection{Lab and demo artefacts}
% HypeBearsVulnerable / Fixed; teaching receiver; Hardhat assertions (a)(b)(c);
% Sepolia Fixed-only page (EOA mint shows no callback window).

\section{Criteria for Assessing Solutions}\label{sec:criteria}
% Table: observability; teaching fidelity; implementation cost;
% new attack surface; classroom demo fitness; alignment with public threats.

\section{Research Methodology}
% What we did (not findings). Literature / public postmortems;
% local cases (HypeBears focus; OMNI/Revest briefly);
% test design (test/hypebears-callback-flow.ts);
% visual-report / animation; Sepolia Fixed-only scope.
% Explicit defensive boundary: no exploit PoC; no public Vulnerable deploy; Toy assets only.

\section{Analysis and Interpretation}
\subsection{Callback results against the criteria}
% Vulnerable vs Fixed (balance 2 vs 1; flag during first callback).
% Table: permission handover vs control-flow handover.

\subsection{Approval thread against the criteria}
% Public-report mechanism alignment; contrast with callback line.

\subsection{Limitations}
% Toy simplification; uneven case depth; EOA demo boundary.
% Point to docs/LIMITATIONS-AND-FUTURE.md.

\section{Conclusions and Recommendations}
% Recycle criteria language; CEI + least privilege;
% Future work from LIMITATIONS-AND-FUTURE.md
% (ERC-1155 always-callback contrast; payment/supply before first safe*;
% Guard depth; fairness/quota framing).

\appendices
\section{Key Lab Artefacts}
% Paths on levi_dev: contracts, tests, CALLBACK-FLOW, SEPOLIA-MINT, LIMITATIONS-AND-FUTURE.

\section{Public Sources for the Approval Hook}
% Check Point / Beosin / CNA / CybersecAsia — public reports only.

\begin{thebibliography}{99}
\bibitem{eip721} W.~Entriken et al., ``EIP-721: Non-Fungible Token Standard,'' Ethereum Improvement Proposals.
\bibitem{eip1155} W.~Radomski et al., ``EIP-1155: Multi Token Standard,'' Ethereum Improvement Proposals.
% Add Check Point, Beosin, CNA, OpenZeppelin docs, lab RESEARCH.md as needed.
\end{thebibliography}

\end{document}
